\begin{afterword}
\summaryhead

The post starts from the prophecy that machines will do all the work. If people get bored, they can play computer games; but games contain effortful tasks, and once programs help with those, one might as well make the games easier, until nothing is left but a screen that says ``YOU WIN'' forever. The author compares this to the Christian Heaven: eternal laziness sounds like good news only to someone who still has to work. A game is synthetic work, and the author suggests that ``getting rid of work'' should become ``removing low-quality work to make way for high-quality work.''

A long-run meaning of life, the post argues, should consist of goals good to pursue and not only good to achieve. Goals such as curing cancer are legitimate, but only their result matters. Passive experiencing is rejected because the author does not want to become ``a happy blob,'' which the author counts as a justification. As with love in an earlier post, a valuable race needs real opponents, one's own effort, and a sentient being to experience it.

\responsehead

Most of the post's positive claims are statements of value, and they are offered as such. The author ``wouldn't actually prefer'' to be a passive experiencer, and the reason given for not becoming ``a happy blob'' is that the author does not want to. As preferences they need no evidence.

Where the post touches facts, the evidence is on its side. Research on flow ties enjoyment to challenges that stretch a person's skills, which is the assumption behind the joke about easier games. In one study, workers reported most of their flow experiences at work rather than in leisure, which fits the proposal to improve work rather than abolish it. Psychologists call activity valued for itself, apart from its product, ``autotelic.''

The main ideas also have well-known precedents, which support the post. In 1930 Keynes foresaw both the problem and the comparison with Heaven: he quotes a charwoman's epitaph about doing ``nothing for ever and ever,'' calls it ``her heaven,'' and names the coming problem of how to occupy leisure. Bernard Suits's \textsc{The Grasshopper} (1978) asked what people would do in a Utopia with no necessary work, and answered that they would play games. Nozick's case against the experience machine has both the ``blob'' and the wish ``to do certain things.''

The post's ``we should look for goals that are good to pursue'' depends on the author's preference, so it applies to readers who share it; the notes on ``Not for the Sake of Happiness (Alone)'' give evidence that many do.

In short: a statement of preferences about a future without necessary work, consistent with research on flow and anticipated by Keynes, Suits and Nozick.
\end{afterword}
